\section{Common observation and change-of-law estimates}\label{sec:observations}

The two geometries use different master widths and nuclear fields, but
share the observation likelihood, the cost of imposing an event, and
the transfer from the observed field to a fresh patch field. We prove
these statements here. In particular, the last transfer compares
expectations under a specified law; it does not identify the two
conditional fields on individual configurations.

\subsection{Likelihoods and conditional versions}

Let $\Psi$ be a normalized antisymmetric state in a finite $n$-electron
sector. Under a probability law $P$, raw positions and spins
$(x,\boldsymbol\sigma)$ have density $|\Psi|^2$. Fix finitely many
widths $\ell_0,\ldots,\ell_{m-1}>0$. Independently of the raw data,
take independent scalar noises $U_{ji,b}$ with density
\[
 \zeta(u)=c_\zeta\exp\!\left(-\frac1{1-u^2}\right)
                         \1_{(-1,1)}(u).
\]
At scale $j$, form the labeled observations
$\widehat Y_{ji,b}=x_{i,b}+\ell_jU_{ji,b}$ and then forget the
particle labels separately in each array. Represent the result by
its lexicographically sorted tuple $Y_j$ in the strict sorted chamber
$\mathcal C\subset(\R^3)^n$. Ties have probability zero. Set
\[
 \mathcal F_j=\sigma(Y_j,\ldots,Y_{m-1}),\qquad
 \mathcal F_m=\{\varnothing,\Omega\}.
\]
Thus the sigma-fields decrease as $j$ increases.

Let $t_z$ be a deterministic measurable width with
$t_z\ge t_{\min}>0$ for the fixed system, and use the smooth radial
packet of Section~\ref{sec:localization} to define
$\mathcal K_z(y)=g_{t_z}(y-z)^2$. Write
$V=V_s+V_c$, where $V_s$ is the deterministic nuclear potential
retained in the observed field. In the molecular application it is
the smoothed nuclear potential. In the atomic application
$V_s=V$ and $V_c=0$. Where relevant put
$\nu_s=-(4\pi)^{-1}\Delta V_s$ distributionally; it is the
smoothed nuclear measure in the first case and the original
point measure in the second.

\begin{lemma}[Conditional densities and fields]\label{lem:obs-data}
There are jointly Borel versions of
\[
 \mu_j(y)=\E_P\!\left[\sum_{i=1}^n\mathcal K_{x_i}(y)
                                      \mid\mathcal F_j\right],
 \qquad H_j(y)=V_s(y)-(K*\mu_j)(y),
\]
with the following properties. At every positive-density data tuple,
$\mu_j\ge0$, $\int\mu_j=n$, and $\mu_j$ is smooth in its spatial
variable. Every spatial derivative has a deterministic bound for
the fixed system. The potential $K*\mu_j$ is $C^1$, with bounded
value and gradient, and
\begin{equation}\label{eq:obs-poisson}
 \Delta H_j=-4\pi\nu_s+4\pi\mu_j
\end{equation}
in distributions. In particular, the always valid regularity
statement is that $H_j-V_s$ is $C^1$ and bounded with bounded
gradient. If $V_s$ and its gradient are bounded and continuous,
these conclusions hold for $H_j$ itself. Conditional averaging
satisfies
\begin{equation}\label{eq:obs-tower}
 \E_P[\mu_j(y)\mid\mathcal F_{j+1}]=\mu_{j+1}(y),\qquad
 \E_P[H_j(y)\mid\mathcal F_{j+1}]=H_{j+1}(y)
\end{equation}
where $V_s(y)$ is finite. The same towers hold for nonnegative
spatial integrals, and for integrable signed observables.
These statements include the trivial data at $j=m$. None of the
fixed-system regularity bounds is asserted to be uniform in the
nuclear system or in $t_{\min}$.
\end{lemma}

\begin{proof}
On $\mathcal C$ the exact likelihood of an array is
\begin{equation}\label{eq:obs-likelihood}
 L_j(x,Y_j)=\sum_{\pi\in S_n}\prod_{i=1}^n\prod_{b=1}^3
 \ell_j^{-1}\zeta\!\left(
           \frac{Y_{j,\pi(i),b}-x_{i,b}}{\ell_j}\right).
\end{equation}
Set it equal to zero outside the chamber. Integration of the
permutation sum over $\mathcal C$ equals integration of the labeled
product over all arrays, and is therefore one. Define
\[
 \Lambda_{\ge j}(x,Y)=\prod_{h=j}^{m-1}L_h(x,Y_h),\qquad
 p_{\ge j}(Y)=\sum_{\boldsymbol\sigma}\int
       |\Psi(x,\boldsymbol\sigma)|^2\Lambda_{\ge j}(x,Y)\dd x,
\]
with the empty product and $p_{\ge m}$ both equal to one. At a tuple
with $p_{\ge j}>0$, use the probability density
$|\Psi|^2\Lambda_{\ge j}/p_{\ge j}$ as the conditional raw law.
Equivalently, the conditional mean of a nonnegative Borel raw
observable $G$ is the ratio
\[
 \frac{\displaystyle\sum_{\boldsymbol\sigma}\int
       G(x,\boldsymbol\sigma)|\Psi(x,\boldsymbol\sigma)|^2
                      \Lambda_{\ge j}(x,Y)\dd x}
      {p_{\ge j}(Y)}.
\]
On zero-density tuples use the unconditional raw law. All these
versions are jointly Borel, including their dependence on an
additional spatial variable. At every parent tuple with
$p_{\ge j+1}>0$, the conditional density of the next array is
$p_{\ge j}/p_{\ge j+1}$. Tonelli's theorem and $\int L_j\dd Y_j=1$
show that it integrates to one. Substitution of the ratios gives
the tower property, also after nonnegative spatial integration;
positive and negative parts give the integrable signed case.

For each multi-index $\alpha$,
$|\partial_y^\alpha\mathcal K_z(y)|
 \le C_\alpha t_{\min}^{-3-|\alpha|}$ uniformly in $z,y$.
Differentiation under the conditional integral therefore gives all
claimed derivatives and bounds. Tonelli gives the mass $n$.
For the potential and its gradient, split the kernels $|u|^{-1}$
and $|u|^{-2}$ at unit distance. The near parts are integrable and
are controlled by $\|\mu_j\|_\infty$; the far parts are controlled
by $n$. Splitting again at a small radius, then using dominated
convergence on its complement, proves continuity of the potential
and gradient. Finally $-\Delta K=4\pi\delta_0$ gives
\eqref{eq:obs-poisson}. In the atomic case this argument does not
remove the nuclear singularity of $V_s$.
\end{proof}

\subsection{Dimension-free energy cost}

\begin{lemma}[Energy cost of a data event]\label{lem:obs-tilt}
Suppose $\Psi$ is an eigenstate of its sector Hamiltonian with
energy $e_\Psi$, where $0\le e_\Psi-E\le D_0$. In particular,
$e_\Psi=E_n$ for a sector ground state. Suppose further that
\[
 \sum_{h=j}^{m-1}\ell_h^{-2}\le C_\ell\ell_j^{-2}
\]
with a fixed numerical $C_\ell$. This holds for the dyadic scales
$\ell_h=r_h^{1.01}$, $r_{h+1}=2r_h$, with
$C_\ell=(1-2^{-2.02})^{-1}$. If $A\in\mathcal F_j$ has probability
$p=P(A)>0$, put $p_A(x)=P(A\mid x)$ and
\[
 \Psi_A=\sqrt{p_A(x)/p}\,\Psi.
\]
Then $\Psi_A$ is a normalized antisymmetric form-domain vector,
its raw position-and-spin law is the raw marginal of $P$ conditioned
on $A$, and
\begin{equation}\label{eq:obs-tilt}
 q_n[\Psi_A]\le e_\Psi+C\ell_j^{-2}\log^5(e/p)
       \le E+D_0+C\ell_j^{-2}\log^5(e/p).
\end{equation}
The constant depends only on $C_\ell$ and the chosen noise law,
not on $n$, the number of arrays, or the nuclear system. The same
assertion holds for an event of one unordered observation at any
width $\ell$, using $\ell$ in place of $\ell_j$. No upper bound
on the particle number in terms of the nuclear charge is needed.
\end{lemma}

\begin{proof}
In labeled observation coordinates place the sorting map inside
the indicator of $A$. The translated product density is continuously
differentiable in $L^1$, since $\zeta$ is smooth and compactly
supported. Hence $p_A$ is continuously differentiable even for a
Borel event. The scalar translation score is
\[
 S(u)=-\frac{\zeta'(u)}{\zeta(u)}
       =\frac{2u}{(1-u^2)^2},\qquad |u|<1.
\]
Its law is symmetric and centered. It satisfies
$\|S(U)\|_q\le Cq^2$ for $q\ge2$: indeed,
$T=(1-U^2)^{-1}$ has a finite exponential moment of every order
less than one, so integration of its exponential tail gives
$\|T\|_q\le Cq$, while $|S(U)|\le2T^2$.

For a unit vector $v\in\R^{3n}$ the directional derivative of the
likelihood pairs the event indicator with
\[
 Q_v=\sum_{h=j}^{m-1}\sum_{i=1}^n\sum_{b=1}^3
                    \ell_h^{-1}v_{i,b}S(U_{hi,b}).
\]
Multiplication by independent random signs preserves the joint law
of the independent symmetric scores. For deterministic coefficients,
\[
 \prod_\alpha\cosh(sc_\alpha)
       \le\exp\!\left(\frac{s^2}{2}\sum_\alpha c_\alpha^2\right)
\]
gives the Gaussian tail estimate and consequently
$\|\sum_\alpha\epsilon_\alpha c_\alpha\|_q
 \le C\sqrt q(\sum_\alpha c_\alpha^2)^{1/2}$.
Conditioning on the scores and then using the triangle inequality
in $L^{q/2}$ yields
\[
 \|Q_v\|_q\le Cq^{5/2}
       \left(\sum_{h,i,b}\ell_h^{-2}v_{i,b}^2\right)^{1/2}
       \le Cq^{5/2}\ell_j^{-1}.
\]
The identity $\sum_{i,b}v_{i,b}^2=1$ and the geometric square-sum
bound are precisely what avoid factors depending on dimension
or on the number of observations.

H\"older's inequality at fixed $x$ gives
$|D_vp_A(x)|\le Cq^{5/2}\ell_j^{-1}p_A(x)^{1-1/q}$.
For $p_A(x)>0$ choose $q=\max\{2,\log(e/p_A(x))\}$ and then take
the supremum over unit vectors. At a zero the gradient of the
nonnegative differentiable function $p_A$ vanishes. Therefore
\begin{equation}\label{eq:obs-score-gradient}
 |\nabla p_A|\le C\ell_j^{-1}p_A\log^{5/2}(e/p_A),
\end{equation}
with right side zero at $p_A=0$.

Regularize the square root by $\sqrt{p_A+\eta}$. Since
$\sqrt u\log^{5/2}(e/u)$ is bounded on $[0,1]$, these functions
have a common Lipschitz bound. Their uniform limit $\sqrt{p_A}$ is
Lipschitz. Its gradient vanishes almost everywhere on its zero set,
and the ordinary chain rule applies elsewhere. Thus
\[
 \left|\nabla\sqrt{p_A/p}\right|^2
       \le C\ell_j^{-2}\frac{p_A}{p}\log^5(e/p_A).
\]
The permutation sum in \eqref{eq:obs-likelihood} shows that $p_A$
is invariant under raw particle permutations and is spin independent.
The bounded Lipschitz multiplier therefore preserves antisymmetry
and the form domain. Its squared norm and its raw law are as claimed
by the definition of $p_A$.

Lemma~\ref{lem:sector-multiplier}, applied with eigenvalue $e_\Psi$,
now gives
\[
 q_n[\Psi_A]-e_\Psi
    \le C\ell_j^{-2}\E_P[\log^5(e/p_A(x))\mid A].
\]
Under this conditional law,
\[
 P\{p_A(x)<e^{-u}\mid A\}
   =p^{-1}\E_P[p_A\1_{\{p_A<e^{-u}\}}]
   \le\min\{1,e^{-u}/p\}\qquad(u\ge0).
\]
Integration of this tail, split at $\log(e/p)$, bounds the fifth
moment by $C\log^5(e/p)$ and proves \eqref{eq:obs-tilt}.
The proof for a single array is identical with the square-sum
constant equal to one.
\end{proof}

\subsection{Transfer to a fresh patch}

The following quantitative formulation isolates exactly the local
screening estimates needed by both applications. It allows either
a patch containing singular nuclei or a patch containing none.
The fine density below is the actual sum of packets around retained
raw particles, not its conditional average.

\begin{lemma}[Averaged field transfer]\label{lem:obs-transfer}
Fix a point $y$ where $V_s(y)$ is finite and a length $a>0$.
Let $A\in\mathcal F_j$ have positive $P$-probability, and let $Q$
be $P$ conditioned on $A$, extended by fresh localization labels.
Its raw marginal is the law of $\Psi_A$ from
Lemma~\ref{lem:obs-tilt}. Let $X$ denote the fresh patch data,
and let $\rho_X^c$ be the conditional density of core particles
under this tilted localization law. Partition the remaining raw
particles into retained and deleted particles. Let
\[
 \sigma=\sum_{i\ {\rm retained}}\vartheta_b(\cdot-x_i),\qquad
 h_X=V_s-K*\rho_X^c-K*\rho_X,\qquad W_X=V_c+h_X,
\]
where $\vartheta_b$ is the radial probability kernel of
Lemma~\ref{lem:loc-cubes}, supported in $B(0,\sqrt3b)$, and
$\rho_X\ge0$ is a measurably chosen patch comparison density.
Assume $\rho_X$ and
$\sigma$ belong to $L^{5/3}$ and have compact support for almost
every patch datum. Put $t=t_y$, $\tau=t/a$, and $\beta=b/a$,
and suppose $0<t,b\le a/10$. Assume the following estimates, with
fixed constants:
\begin{enumerate}
\item The master width $z\mapsto t_z$ is Lipschitz with constant
at most $1/2$. The raw one-particle density $\rho_{\Psi_A}$ satisfies
\[
 \int_{B(y,4a)}\rho_{\Psi_A}^{5/3}\le Ca^{-7}.
\]
\item Deleted particles have $|x_i-y|\ge ca$, and their count obeys
$\E_Q n_{\rm del}\le Ca^{-3}\sqrt\beta$.
\item For some $\varepsilon_a\ge0$,
\[
 \E_Q D(\rho_X-\sigma)\le Ca^{-7}\varepsilon_a^2.
\]
\item On $B(y,t)$ either the source-free bound
$\rho_X\le Ca^{-6}$ holds, or
\[
 \rho_X=k(W_X)_+^{3/2},\qquad
 W_X\le V_{\rm loc}+Ca^{-4},\qquad
 V_{\rm loc}(z)=\sum_\alpha\frac{q_\alpha}{|z-R_\alpha|},
 \qquad\sum_\alpha q_\alpha\le Ca^{-3},
\]
with $q_\alpha\ge0$. The latter alternative permits arbitrary
locations of the local nuclei.
\end{enumerate}
Then the two expectation towers are
\begin{equation}\label{eq:obs-two-towers}
 \E_Q(K*\mu_j)(y)=\E_Q P^{\rm ms}(y),\qquad
 \E_Q(K*\rho_X^c)(y)=\E_Q P^c(y),
\end{equation}
where $P^{\rm ms}$ is the potential of the raw master packets and
$P^c$ is the raw core potential. Moreover,
\begin{equation}\label{eq:obs-transfer}
 \begin{split}
 |\E_Q(H_j-h_X)(y)|\le Ca^{-4}\big[&
    \tau^{1/5}+\beta^{1/5}+\sqrt\beta
      +\varepsilon_a\tau^{-1/2}\\
    &+R(\tau)+(\tau+\sqrt3\beta)^{1/5}\big],
 \end{split}
\end{equation}
where $R(\tau)=\tau^2$ in the source-free alternative and
$R(\tau)=\tau^2+\tau^{1/2}$ in the nuclear alternative.
Thus the common upper bound $R(\tau)\le C\tau^{1/2}$ is always
available. All constants are independent of the number of particles;
they depend only on the constants in the displayed hypotheses.
\end{lemma}

\begin{proof}
Write
\[
 P^{\rm raw}(y)=\sum_iK(y-x_i),\qquad
 P^{\rm ms}(y)=\sum_i(K*\mathcal K_{x_i})(y),\qquad
 P^{\rm raw}=P^c+P^{\rm ret}+P^{\rm del}.
\]
The first identity in \eqref{eq:obs-two-towers} follows by multiplying
the original $P$-conditional expectation defining $\mu_j$ by
$\1_A$, using $A\in\mathcal F_j$, and dividing by $P(A)$.
The fresh labels do not alter this expectation. For the second
identity, condition the raw-plus-label marginal of $Q$ on $X$.
That marginal is exactly the localization law of $\Psi_A$; the
original observations can be integrated out. Tonelli justifies
both identities initially for nonnegative potentials. In
particular the first tower is under the original observation
law, while the second uses the fresh tilted patch law.

Subtracting the definitions of $H_j$ and $h_X$, and then using
these towers, gives the exact algebraic identity
\begin{equation}\label{eq:obs-transfer-decomposition}
 \begin{split}
 \E_Q(H_j-h_X)(y)=\E_Q\big[&
 (P^{\rm raw}-P^{\rm ms})-P^{\rm del}\\
 &-(P^{\rm ret}-K*\sigma)+K*(\rho_X-\sigma)\big](y).
 \end{split}
\end{equation}
The estimates below also establish finiteness of all terms, so this
subtraction involves no difference of infinite expectations.

For $0<u\le4a$, the local density bound and H\"older's inequality
give
\begin{equation}\label{eq:obs-near-coulomb}
 \begin{split}
 \E_Q\sum_{|x_i-y|<u}K(y-x_i)
 &\le\left(\int_{B(y,4a)}\rho_{\Psi_A}^{5/3}\right)^{3/5}
       \left(\int_{|z|<u}|z|^{-5/2}\dd z\right)^{2/5}\\
 &\le Ca^{-4}(u/a)^{1/5}.
 \end{split}
\end{equation}
The complementary raw potential is bounded by $n/(4a)$.
By radial smearing, $P^{\rm raw}-P^{\rm ms}\ge0$ and only centers
with $|x_i-y|\le t_{x_i}$ can contribute. Lipschitz continuity of
the width implies $|x_i-y|\le t_y+|x_i-y|/2$, hence these centers
lie in $B(y,2t)$. Equation~\eqref{eq:obs-near-coulomb} bounds its
mean by $Ca^{-4}\tau^{1/5}$. The same Newton comparison for each
fine packet gives
\[
 0\le\E_Q(P^{\rm ret}-K*\sigma)(y)
             \le Ca^{-4}\beta^{1/5}.
\]
The distance and count assumptions for the deleted particles give
$\E_Q P^{\rm del}(y)\le Ca^{-4}\sqrt\beta$.

For the signed density difference take
$v_t(z)=\min\{|y-z|^{-1},t^{-1}\}$. This is in $\dot H^1$ and
\[
 \|\nabla v_t\|_2^2=4\pi\int_t^\infty r^{-2}\dd r=4\pi/t.
\]
Lemma~\ref{lem:loc-coulomb} and Cauchy--Schwarz in probability give
\[
 \E_Q\left|\int v_t(\rho_X-\sigma)\right|
   \le Ct^{-1/2}\sqrt{\E_QD(\rho_X-\sigma)}
   \le Ca^{-4}\varepsilon_a\tau^{-1/2}.
\]
It remains to control the two nonnegative remainders of the
truncation inside $B(y,t)$. In the source-free case,
\[
 \int_{B(y,t)}K(y-z)\rho_X(z)\dd z
                         \le Ca^{-6}t^2=Ca^{-4}\tau^2.
\]
For the nuclear alternative, weighted convexity gives, with
$Q_{\rm loc}=\sum_\alpha q_\alpha$,
\[
 V_{\rm loc}^{3/2}
   \le Q_{\rm loc}^{1/2}\sum_\alpha
                           q_\alpha|z-R_\alpha|^{-3/2}.
\]
The assertion is read as zero if $Q_{\rm loc}=0$. Uniformly in
$R\in\R^3$,
\[
 \int_{B(y,t)}\frac{\dd z}{|y-z|\,|z-R|^{3/2}}
                                                  \le Ct^{1/2}.
\]
For completeness, apply H\"older to the second factor with any
exponent $q\in(3/2,2)$ and to the first with conjugate exponent
$q'\in(2,3)$. Integrals of these radial decreasing powers over a
ball are no larger than their centered-ball integrals; the powers
of $t$ combine to $t^{1/2}$. Since
$\rho_X\le C(a^{-6}+V_{\rm loc}^{3/2})$, its remainder is at most
\[
 C\big(a^{-6}t^2+Q_{\rm loc}^{3/2}t^{1/2}\big)
                    \le Ca^{-4}(\tau^2+\tau^{1/2}).
\]
This proves the asserted nuclear bound even when a nucleus lies at
$y$; only $V_s(y)$, not the unsmoothed $V(y)$, was required finite.

For $\sigma$, only packets centered within $t+\sqrt3b$ of $y$
meet the truncation ball. Extend each such packet's kernel integral
to all space and use radial smearing to bound it by its raw point
potential. Equation~\eqref{eq:obs-near-coulomb} bounds the mean
remainder by $Ca^{-4}(\tau+\sqrt3\beta)^{1/5}$. The truncated
$\sigma$ potential is integrable, being at most $n/t$. Together
with Coulomb duality and the just obtained remainder bound, this
also proves integrability of the full $\rho_X$ potential. The
raw and core potentials were already integrable by
\eqref{eq:obs-near-coulomb}. Hence all the preceding subtractions
are valid. Adding the four estimates in
\eqref{eq:obs-transfer-decomposition} proves \eqref{eq:obs-transfer}.
\end{proof}

In both later applications the screening estimates supply
$\varepsilon_a=a^{.01}$ and $\beta=a^{.2}$. The master geometry
supplies $\tau\ge ca^w$, with $w=10^{-5}$, and an upper bound on
$\tau$ tending to zero with the outer small parameter. Consequently
\[
 a^{.01}\tau^{-1/2}\le Ca^{.01-w/2},\qquad
 .01-w/2=.009995>0,
\]
and every term in the bracket in \eqref{eq:obs-transfer} tends to
zero uniformly in the participating bands. Thus the transfer is
$o(a^{-4})$. The specific lower probability thresholds for the
events are used in those applications to verify the screening
hypotheses through \eqref{eq:obs-tilt}; no probability threshold is
needed for the transfer identity itself.
